\section[第8章\q 克服困难]{第8章\\克服困难}

人们通过学习掌握了财务知识，但在通向财务自由之路上仍面临着许多障碍。资产项可以产生大量的现金流，使人们过上梦想中的生活，而不必整天为了支付账单而忙碌。但掌握财务知识的人有时候还是不能积累丰厚的资产项，其主要原因有5个：

{\kaishu
1. 恐惧心理。

2. 愤世嫉俗。

3. 懒惰。

4. 不良习惯。

5. 自负}

原因之一：对可能亏钱的恐惧心理。

我从未遇见喜欢亏钱的人，但在我的一生中，从未遇见一位没亏过钱的富人，却遇见许多从未亏掉一毛钱的穷人——我是说在投资活动中。

对亏钱的害怕是确实存在的，每个人都有这种恐惧心理，包括富人。但恐惧本身并不是问题之所在，关键在于你如何处理恐惧心理，如何处理亏钱的问题。处理失败的不同方式造成了人们生活的差异，不仅是对金钱，对生活中的任何事情都是如此。富人和穷人之间的主要差别在于他们处理恐惧心理的方式不同。

感到恐惧是正常的，在涉及金钱时表现出胆怯也是正常的，即便如此你仍然有机会成为富人。我们在某些方面是英雄，而在另外一些方面却可能是懦夫。我朋友的妻子是一位急诊室护士，当她看到流着血的病人时，会冲上去救治，可是当我提到投资时，她却逃跑了。不过当我看到血时，我不会跑，我会直接晕倒。

我的富爸爸理解人们对金钱的恐惧症。“一些人害怕蛇，一些人害怕失去金钱，这都是恐惧症。”他说。因此，他有个克服这种恐惧的小窍门：“如果你讨厌冒险，担心会亏钱，就早点动手积累资产吧。”

这也是银行建议你在年轻时就养成储蓄的习惯的原因。如果在年轻时就开始积累，你就更容易致富。我不会在这里详细讨论这个问题，但是从20岁开始储蓄与从30岁开始储蓄的人之间的确有巨大的差异。

有人说这世界的奇迹之一就是复利。据说，购买曼哈顿岛是有史以来最廉价的交易之一\footnote{1626年，荷兰人用价值24美元的物品从印第安人手中买下了曼哈顿岛，并取名为新阿姆斯特丹。1664年曼哈顿岛被英国人占领，改名为纽约。}，它是用价值24美元的廉价小玩艺儿买下来的。然而，如果将那24美元用于投资，以8％的年利率计算，到1995年这24美元就会变成28万亿美元。如果以1995年的房地产价格计算，用这笔钱不仅可以买回曼哈顿岛，还有余钱几乎能买下洛杉矶。

我的邻居为一家大型电脑公司工作，他在那儿干了25年。再过5年，他就会离开那家公司，按401(k)计划他将得到400万美元。这些钱大部分被投资于高增长的共同基金，他会将其转换成债券和国债。他离开公司时只有55岁，每年却能获得超过30万美元的现金流入，这比他的工资收入还要高。所以，如果你害怕亏钱或者讨厌冒险，你至少可以像他一样。但是，你必须及早行动并且制定一个完善的退休计划，此外你还要聘请一位信得过的财务规划师，在你作出任何投资决定之前，他能给你指导。

可是，如果你没有足够的时间或是希望早点退休，又该怎么办呢？

你怎样应付亏钱的恐惧心理呢？

我的穷爸爸在这方面没有任何作为。他一味回避这个问题，拒绝进行讨论。

我的富爸爸恰恰相反，他建议我要像得克萨斯人那样思考。“我喜欢得州和得克萨斯人，”他曾说，“在得州，什么东西都大气。如果得克萨斯人赢了，他们就会赢很多；如果他们输了，也会输得令人吃惊。”

“难道他们喜欢失败吗？”我问道。

“我不是这个意思，没有人喜欢失败。有句话说得好，‘没有快乐的失败者’，”富爸爸说，“我要说的是得克萨斯人对于风险、收益和失败的态度。那是他们驾驭生活的方式，他们活得很大气，不像这儿的大部分人在碰到金钱问题时，就变得像斜齿鳊一样。斜齿鳊在被光照到时会非常害怕，而这种人在杂货店店员少找两毛五分钱时，便会抱怨不停。”

富爸爸接着解释：“我最喜欢的是得克萨斯式的生活态度，得克萨斯人赢了会骄傲，输了也不忘吹牛。他们有一句谚语，‘如果你要破产了，那就索性一分不剩’。我们身边的大部分人都非常害怕亏钱，而得克萨斯人却不愿意让你认为他们仅仅因为一幢复式公寓而破产。”

富爸爸经常跟我和迈克说，大部分人在财务上不成功的最大原因是他们的做法过于安全。“因为太害怕失败，所以才会失败。”他常这么说。

对此前美国职业橄榄球联赛的四分卫弗朗·塔肯顿还有另一种说法：“胜利意味着不害怕失败。”

在我的生活中，我注意到失败与成功相伴相随。在我学会骑自行车之前，曾经跌倒过许多次。我从未遇见不曾打丢一球的高尔夫球选手，也从未见过不曾伤心过的恋人，更未见过从不亏钱的富人。

因此，对大多数人来说，他们在财务上不成功是因为对他们而言亏钱所造成的痛苦远远大于致富所带来的乐趣。得克萨斯的另一句谚语讲道：“人人都想上天堂，却没有人想死。”大部分人都梦想发财，但却害怕亏钱，所以他们永远也进不了财富的天堂。

富爸爸常常给我和迈克讲他到得州旅行的故事。“如果你真想学习如何面对风险、损失和失败，就去圣安东尼奥市走访阿拉莫。”阿拉莫的故事讲的是，一群勇敢的人在知道毫无胜算的情况下依然选择战斗，他们宁可战死也不愿投降。这是一个值得学习的、振奋人心的故事，然而，这的确是一次悲壮的军事失败。他们最后不得不用枪托进行肉搏。

你完全可以称之为失败，他们最终失守阿拉莫。你想知道得克萨斯人是如何面对失败的吗？他们高声呼喊：“记住阿拉莫！”

我和迈克多次听到这个故事。每次在做大生意之前或者感到不安的时候，富爸爸就会给我们讲这个故事；当他把一切仔细安排好或者一件事情尘埃落定的时候，他也会讲这个故事；当他担心犯错误或害怕亏钱时，他还会讲这个故事。这个故事给了富爸爸力量，因为它总在提醒他，他可以将财务上的损失变成赢利。富爸爸知道失败只会使他更强大，更精明。他并不愿意亏钱，但他清楚自己是什么样的人，知道该怎样去面对损失。他接受了损失，然后将它变成赢利，这是他成为赢家而别人沦为失败者的原因所在；同时这也是当别人退出时，他依然有勇气去冲过终点线的原因。“这就是我这么喜欢得克萨斯人的原因。他们接受失败的现实并把它转变成通向成功道路上的一个个插曲。”

今天我对富爸爸的这番话有了更深的体会：“得克萨斯人并不掩饰失败，而是愈挫愈奋，他们接受失败的现实并将失败转化为动力。失败激励得克萨斯人走向成功，而这个道理并不仅仅适用于得克萨斯人，它适用于所有的成功人士。”

就像我前面说过的：从自行车上摔下来是学习骑车的一部分。我还记得正是摔下来才使我更坚定地要学会骑车。同样的，我也说过世界上没有从未打失过一球的高尔夫球选手。作为一位职业高尔夫球高手，打失一个球或输掉一场比赛只会激励他做得更好，练得更努力，学更多的东西，是失败使他们更加优秀。对于胜利者，失败会激励他们；对于失败者，失败则会击垮他们。

用洛克菲勒的话来说，就是“我志在将每一次灾难转化成机会”。

作为一个日裔美国人，我也可以说这样的话。许多人说珍珠港事件是美国人的失误，我却认为是日本人的最大失误。在电影《虎！虎！虎！》中，一个日本海军上将悲哀地对他的兴高采烈的部下说：“我担心我们摇醒了一个沉睡的巨人。”果然“记住珍珠港”成为一句具有巨大感召力的口号，它把美国历史上最大的失败之一变成了取得胜利的原因，这次巨大的失败给了美国力量，此后美国很快就崛起为一个世界强国。

失败会激励胜利者，击垮失败者。这是胜利者最大的秘密，也是失败者所不知道的秘密。胜利者最大的秘密是失败能够激励他们取胜，所以他们不怕失败。让我们重复一遍弗朗·塔肯顿的话：“胜利意味着不害怕失败。”像弗朗·塔肯顿这样的人不害怕失败，因为他们了解自己。他们和所有人一样不喜欢失败，但他们知道，失败会激励他们做得更好。

要知道不喜欢失败和害怕失败之间有着巨大的差别，大部分人因为害怕失败而失败，他们甚至会因一幢复式公寓而破产。在财务上他们做得过于安全、规模太小，他们会买大房子、大轿车，却不进行大的投资。

90％以上的美国人财务困难的主要原因就在于，他们是为了避免损失而理财，而不是为了赢利而理财。

他们会选择去找财务规划师、会计师和股票经纪人等，购买一个安全的投资组合。他们中的大部人将大量现金以存款单、低收益债券、可以在共同基金内部买卖的共同基金，以及一点股票的形式进行投资。这是一个安全而合理的投资组合，却并不是一个赢利的投资组合。说到底这种投资组合只是人们为避免损失而采取的方式。

对于超过70％的人来说，这可能还算是一个较好的投资组合。毕竟，一个安全的投资组合总比什么都没有强得多。一个安全的投资组合对于注重安全的人来说是很合适的，但是，安全地、平衡地投资于一个投资组合，却不是一个成功的投资者应有的投资方式。如果你的资金很少而又想致富，你必须首先集中于一点，而不是追求平衡。那些成功的投资者，在最初肯定不是追求平衡的，追求平衡的人只会在原地踏步。

要进步，你必须先做到不平衡，并注意怎样才能使自己不断进步。

爱迪生不追求平衡，他集中精力于某样东西；比尔·盖茨也不追求平衡；唐纳德·川普把注意力仅放在一点上；乔治·巴顿从不把坦克部署在很长的战线上，而是集中起来攻击德国防线上最薄弱的地方，与此相反，法国人构筑了漫长的马其诺防线，其结局众所周知。

如果你有致富的愿望，就必须集中精力。把你大部分的鸡蛋放在较少的篮子里，别像穷人和中产阶级那样：把很少的鸡蛋放在许多篮子里。

如果你不愿失败，那就安全地投资；如果损失会使你元气大伤，那就稳妥一点，去做一个“平衡的”投资。你要是已经超过了25岁并害怕冒险，那就不要改变自己的投资方式。以安全的方式投资，但要尽早起步，要早点开始积累你的“鸡蛋”，因为以这种方式积累需要大量的时间。

然而，假如你梦想获得自由——从“老鼠赛跑”的游戏中解脱出来，你应该问自己的第一个问题是：“我该如何去面对失败？”如果失败能激励你去争取胜利，也许你就应该去争取每一次投资机会——但仅仅是也许而已；如果失败会使你损失惨重，或者使你烦躁不安——你就像一个愣头青一样，只要不如意就打电话找律师帮你提起诉讼——那你最好还是做稳妥一点的投资。继续你的日常工作，或者购买一些国债和共同基金，但是要记住，这些金融工具也同样有风险，虽然它们较为安全一些。

我说了这么多，还提到了得克萨斯和弗朗·塔肯顿，只是想说明积累资产项其实非常容易，这就像一场低智商的游戏，不需要多高的教育，五年级的数学水平就够了。然而，用资产进行投资却是一种高智商的游戏，需要胆量、耐心和对待失败的良好态度。失败者回避失败，而失败本来是可以使失败者转变为成功者的。所以一定要“记住阿拉莫”。

\subsection*{原因之二：愤世嫉俗的心理}


“天要塌下来了，天要塌下来了。”很多人都知道“小鸡的故事”，小鸡总是围着谷仓转，警告大家即将到来的厄运。我们知道有的人也爱这么做，其实我们每个人的心里也都有一只“小鸡”。

就像我前面指出的，愤世嫉俗的人就像“小鸡”一样，每当心里害怕、怀疑的时候，他们就会像“小鸡”一样杞人忧天。

我们会对自己产生怀疑：“我不够精明”、“我不够好”、“谁谁都比我强”等等，怀疑常常使我们寸步难行。我们总是会问自己有关“如果”的问题，例如：“如果经济恰好在我投资之后崩溃怎么办”，“如果我无法偿还借款怎么办”，或者“如果事情不像我想的那样该怎么办”。有时我们的朋友或是亲近的人会不由自主地提醒我们有某些缺点，他们常常会说，“你怎么就认为你可以这么做”，或者说“如果这是一个好点子，那其他人怎么不做呢”，或者是“这起不了什么作用，你根本不知道自己在说什么”。这些怀疑的话如此刺耳，以至于我们无法将自己的计划付诸行动，害怕的感觉在心中蔓延，有时我们甚至因此夜不能寐。我们无法前进，我们只能守着那些安稳的东西，看着机会从我们身边溜走。我们眼睁睁地看着时光流逝，心中的结使我们碌碌无为。在生活中我们或多或少都会经历这样的状态。

彼得·林奇来自“忠诚马吉兰”共同基金，他把天要塌下来的警告比做“噪音”，而我们都听过这样的“噪音”。

“噪音”既有我们头脑中自己产生的，也有来自我们外部的，例如：朋友、家人、同事和新闻媒体。林奇回忆，在20世纪50年代，新闻媒体中充斥着核战争威胁论，人们开始修筑防辐射掩体，储存食物和水。如果他们将资金明智地投在市场上，而不是用来修建防辐射掩体，他们今天可能已经实现财务自由了。

几年前当洛杉矶爆发骚乱时，全国枪支的销售额都上升了。在华盛顿州，有个人因为吃了汉堡包中的生肉死了，于是亚利桑那州的卫生部门命令餐馆将牛肉完全煮熟。一家药品公司在一家全国性的电视台播放了一则广告，说人们患上流感。广告是在2月份播出的。患感冒的人数和感冒药的销售额一起增长。

大部分人之所以贫穷，是因为在他们想要投资的时候，周围到处是跑来跑去的“小鸡”，叫嚷着“天要塌下来了，天要塌下来了”。“小鸡”们的说法很有影响力，因为我们每个人的心中也都有一只“小鸡”。因此，我们常常需要极大的勇气，不让谣言和杞人忧天的怀疑加剧我们的恐惧和疑虑。

1992年，我的一个叫理查德的朋友从波士顿来到凤凰城探访我和我妻子，他对我们买卖股票和经营房地产非常感兴趣。当时凤凰城的房地产价格非常低，我们花了两天时间，向他介绍那些在我们看来是获取现金流和资本利得的极好的机会。

我和我妻子并不是房地产中介，我们只是投资者。在调查了一处位于旅游胜地的单元房的情况后，我们给一家房地产中介打电话，当天下午他就将这套单元房卖给了我的朋友。一套两居室的住宅售价仅为4.2万美元，而类似的单元房要卖到6.5万美元。他找到了一笔便宜的买卖，很兴奋地买下来，然后回波士顿去了。

两周后，那家房地产中介打电话给我，说我的朋友反悔了。我立即给他打电话，想弄清楚原因。他说他和他的邻居说了这件事，邻居说这是一笔糟糕的交易，他给的价格太高了。

我问理查德，他的邻居是不是一位投资家，理查德说不是。当我问他为什么要听邻居的话时，他略带防备地说他只是想再观望一阵。

到了1994年，凤凰城的房地产市场开始回暖。那套小单元房的租金每月可达1000美元，冬天最高时达到2500美元。1995年这套单元房的价格为9.5万美元。理查德当时需要投入的全部资金仅为首付5000美元，那样他就可以脱离“老鼠赛跑”的游戏了。而今天，他仍然一事无成。凤凰城的房地产仍有廉价交易的机会，不过你要更努力地去寻找。

理查德的反悔并未让我惊讶，这被称为买家违约。这种心理影响着我们所有人。当我们反悔时，“小鸡”得逞了，而我们实现财务自由的机会却丧失了。

在另一个例子中，我通过持有一小部分拥有税收留置权\footnote{所谓留置权是指债权人按照合同约定占有债务人财产，在与该物有牵连关系的债权未受清偿前，有留置该财产，并就该财产有限受偿的权利。负有纳税义务的纳税人如有欠缴税款}的资产，来替代大额存单投资。我每年能拿到16％的利息，大大高于银行5％的利率。这种权利受到国家房地产法和州法律的保障，而且这种保障比大多数银行更可靠。这种投资方式十分安全，只是缺乏流动性，所以我把它们看做2～10年期限的大额存单。然而几乎每次当我告诉某人（特别是当他拥有大额存单投资时），我以这种方式持有资金，他们都会告诉我这样做太冒险。他们还会给我分析为什么我不应该那样做，但当我问他们是从哪儿得到这种依据时，他们就会说是来自朋友或投资杂志。他们从未这样投资，却总是劝别人不要这样做。我能接受的最低收益率为16％，可那些疑虑重重的人却愿意接受5％的收益率。怀疑的代价真是太高昂了。

我认为疑虑和愤世嫉俗的心态使大多数人安于贫困。生活等着你去致富，可就是这些疑虑使人们无法摆脱贫穷。正如我曾说的，摆脱“老鼠赛跑”的生活在技术上讲十分容易，不需要接受太多的教育，可那些疑虑使大多数人寸步难行。

“愤世者永远不会成功。”富爸爸说。“未经证实的怀疑和恐惧会使人们成为愤世嫉俗者。愤世者抱怨现实，而成功者分析现实。”这是富爸爸最喜欢说的另一句话。富爸爸解释，抱怨蒙蔽人的头脑，而分析使人心明眼亮。通过分析能使成功者看到那些愤世者无法看到的东西，发现被其他人都忽视的机会，而发现机会的能力正是取得成功的关键。

对每一个寻求财务独立或财务自由的人来说，房地产都是一个强有力的投资工具。甚至可以说是独一无二的投资工具。然而，每当我提到房地产时，经常有人说：“我不想去修理厕所。”这就是林奇所说的“噪音”，也是富爸爸所说的愤世者的说法。这种人只会批评和抱怨，而不去分析现实。有些人宁可让疑虑和恐惧蒙蔽自己的思想，也不愿睁开眼睛去分析现实。

因此当某人声称“我不想去修理厕所”时，我都想反击说“你凭什么认为我想去”，他们似乎把修理厕所看得比他们想要的东西更重。我在讨论从“老鼠赛跑”中获得自由，他们却只把注意力放在厕所上，这就是使他们生活贫穷的思维模式。他们总是抱怨而不是去分析。

“‘我不想要’是成功的关键。”富爸爸这样说。

正因为我也不想去修理厕所，我才费了很大劲儿找到了一位物业管理者为我修理厕所。找到一位好的物业管理者来为我管理房产，才能够帮助我增加现金流。更重要的是，一位好的物业管理者有助于我去买入更多的房产，因为我不用去考虑修厕所的事了。一位优秀的物业管理者是房地产交易中成功的关键。因此对我来说，寻找一位好的物业管理者比买卖房地产本身更重要。此外，一位好的物业管理者通常会比房地产中介更先打听到大额交易的消息，这就使他们更加重要了。

这就是富爸爸所说的“‘我不想要’是成功的关键”这句话的真义所在。因为我也不想去修理厕所，我才想出购买更多的房地产并将自己从“老鼠赛跑”中尽快解脱出来的办法。那些一直说“我不想去修理厕所”的人总是拒绝使用这个强有力的投资工具，修厕所总是比他们的财务自由重要。

在股票市场上，我也经常听到人们说，“我不想亏钱”。我不知道是什么让他们认为我或其他投资者喜欢亏钱。他们挣不到钱是因为他们选择不亏钱，他们不去分析实际情况，而只是对另一种强有力的投资工具——股票——不予理睬。

1996年12月，我和一位朋友开车经过邻近地区的一个加油站。我的朋友抬头看了看，发现油价上涨了。我的朋友是那种总是忧心忡忡、“小鸡”式的人，对他来说，天似乎总像是要塌下来了，而且像是要压在他的头上。

当我们到家时，他给我列举了一大堆数据，来说明为什么在未来几年油价会上涨。我以前从未读过这些数据，即使我已经拥有了一家营运中的石油公司的大部分股份。根据这些信息，我立即开始寻找并最终找到了一家新成立的价值被低估的石油公司，这家公司正在勘探新油田，发现这家新公司使我的经纪人非常兴奋。后来我买下了它的65％的股份，共1.5万股。

在1997年2月，我还是和这位朋友，开车经过同一个加油站。的确没错，每加仑汽油的价格上涨了约15％，这位“小鸡”式的朋友非常担忧并且不停地抱怨。我笑了，因为在1997年1月，那家小型石油公司找到了石油。自从那次给我分析了那些数据以后，我买下了1.5万股股票，现在每股价格已上涨到3美元以上。如果我朋友的话是正确的，石油价格还会继续上扬。

这些“小鸡”式的人不去分析问题，而是封闭自己的思想。如果大多数人懂得股票市场上“横盘”（预定低点抛售）是什么意思的话，就会有更多的人为了赢利而投资，而不是为了避免损失而投资。“横盘”好比一个计算机指令，当股价开始下跌时自动卖出股票，帮助你将损失最小化、收益最大化。对于那些害怕受到损失的人来说，这是一项极好的工具。

因此，每当我听到人们执迷于“我不想要”，而不去注意他们想要的东西时，我就知道他们脑子里的“噪音”一定很响。“小鸡”控制了他们的思维，叫喊着“天要塌下来了，厕所坏了”。于是，他们避开了自己的“不想要”，却为此付出了巨大的代价——他们可能永远得不到自己想要的。

富爸爸教给我一种看待“小鸡”的方式，“要像桑德斯上校那样去做”。桑德斯上校在66岁时失去了所有的产业，开始靠社会保险金生活，而那点钱根本不够用。于是他走遍全国推销他的炸鸡秘方，在最终有人肯买秘方之前，他被拒绝了1009次。然而通过不懈努力，他在大部分人都退休的年龄成为了千万富翁。“他是一个勇敢、坚韧不拔的人”，富爸爸说的就是肯德基的创始人哈兰·桑德斯上校。

所以，如果你疑虑重重，又有点害怕，不妨像桑德斯上校那样对待自己内心的“小鸡”：油炸它。

\subsection*{原因之三：懒惰}

忙碌的人常常是最懒惰的人。也许我们都听说过某个人努力工作挣钱的故事，他拼命工作希望让妻子儿女生活得更好。他整天待在办公室里，周末还把工作带回家去做。一天，他回到家，却发现人去楼空，他的妻子带着孩子离开了他。他早就知道自己和妻子之间有一些问题，可他却宁愿选择工作，而不去改善双方的关系。这件事让他非常沮丧，所以在工作中的表现也不如以前了，最后失去了这份工作。

我经常遇到那些过分忙于工作而不关心自己的财富的人。还有一些人过分地忙于工作而不照顾自己的身体。使这两种人如此忙碌的原因是一样的，他们把忙碌作为逃避问题的借口。没有人告诉他们这些，但他们心里其实很明白。事实上，如果你去提醒他们，他们往往还会很不高兴。

如果他们不忙着工作或是与孩子在一起，就会忙着看电视、钓鱼、打高尔夫球和购物。总之，他们内心很清楚自己是在逃避一些很重要的事情。这是懒惰最普遍的表现形式，一种通过忙碌掩饰下的懒惰。

那么，怎么才能治疗这种惰性呢？答案就是要“贪婪”一点。

对我们许多人来说，我们是在把贪婪或欲望看做是坏事的环境中成长起来的。“贪婪的人都是坏人。”我妈妈常常这样说。然而，我们的心里都渴望拥有那些美好、新颖和令人兴奋的东西。因此，父母就常常利用负罪感来抑制这种欲望。

“你只考虑自己，难道你不知道还有兄弟姐妹吗？”这是我妈妈最爱说的一句话。“你还想我给你买什么？难道你认为我们是摇钱树吗？你认为钱是从树上掉下来的吗？你知道我们不是有钱人”则是我爸爸最常说的话。

话虽不多，但其中的负罪感深深地影响了我。

另外还有一种父母，他们会从相反的方面利用这种负罪感，他们会这样说：“我省吃俭用给你买这个，因为我小时候从未得到过它。”我有一个邻居可以说是身无分文，但他的车库里却满是他的孩子的玩具，以至于车都停不进去。父母对于那些被宠坏的孩子有求必应，“我不想让他们尝到贫穷的滋味”是他每天都要说的话。他没有为孩子上大学或自己退休留下一点积蓄，可他的孩子却拥有所有的新玩具。最近银行刚给他邮了一张信用卡，就带上孩子去拉斯维加斯玩了。“我这么做全是为了孩子。”临走时他带着自我牺牲的骄傲神情对我说。

富爸爸禁止我们说“我可付不起”这类的话。

在我自己的家里，这可是我经常听到的。但富爸爸要求他的孩子们这样说：“我怎样才能付得起？”他的理由是：“我可付不起”这句话禁锢了你的思想，使你无法进一步思考。“我怎样才能付得起”这句话则开启了你的头脑，迫使你去思考并寻求答案。

但最重要的是，他觉得“我可付不起”是一句谎言，他坚信人的精神能够做到一切。“人类的精神力量非常非常强大”，他常说，“它知道你能做成任何事。”当你的头脑中那个懒惰的思想说“我可付不起”时，两种思想的交锋就开始了。你的精神愤怒了，而你的懒惰的思想却开始为自己辩护。你的精神大叫：“来吧，让我们去健身房锻炼。”而懒惰的思想会说：“可我太累了，我今天工作真的很辛苦。”你的精神会说：“我厌倦了贫穷的生活，让我们脱离这种命运致富吧。”懒惰的思想则会说：“富人很贪婪，而且致富太麻烦了。这不安全，我可能会亏钱。我要努力工作。我还有许多工作要做。看看我今晚必须做的事情，我的老板希望我明早之前干完这些。”

“我可付不起”带来的悲哀和无助感会使人们失望、迟钝以至意志消沉。“我怎样才能付得起”则打开了充满可能性的快乐和梦想之门。”因此，富爸爸并不关心我们想买什么，他只想通过促使我们不断思索“我怎样才能付得起”来创造一种更强有力的思想和更有活力的精神。

所以，他很少给我和迈克买东西，相反，他会问：“你怎样才能买得起这个？”于是包括上大学的学费，都是我们自己挣钱支付的。他希望我们学习的并不是目标本身，而是达到目标的过程。

我认为，今天的问题是成千上万的人对自己的“贪婪”感到内疚，这是他们在少年时代就养成的陈旧思维。他们渴望拥有生活中那些更美好的东西，但大部分人却下意识地调整了自己的心态，并对自己说：“你不能拥有这个”，或是“你可付不起”。

那么，你怎样才能克服懒惰的心理呢？答案是多一点“贪婪”。记得有家电台的频率是WII-FM，它可以被解释成“我还能得到什么”（What'sIn It-For me）。人们需要坐下来问问自己：“如果我身体健康、性感、长相英俊，我还能得到什么？”或者：“如果我不再工作，我的生活会是什么样？”又或：“如果我拥有足够的钱，我会做什么？”没有一点“贪婪”，没有想拥有更好的东西的渴望，就不会前进。世界之所以发展是因为我们都渴望生活得更好，新发明的诞生是因为我们渴望更好的东西，我们努力学习也是因为我们想要了解更好的东西。因此，每当你发现自己在逃避你内心清楚应该去做的事情时，就应该问问自己：“我还能得到什么？”稍稍“贪婪”一点，这是治愈懒惰的灵丹妙药。

当然，就像任何事情都要有“度”一样，过于贪婪就不好了。但要记住迈克尔·道格拉斯在电影《华尔街》中说的话：“欲望是好事。”富爸爸用另一种方式进行解释：“负罪感比欲望要糟，因为负罪感抢走了灵魂。”对我来说，埃莉诺·罗斯福说的最好：“做你内心认为正确的事情因为你不管怎么做总会受到批评。如果你做，会受到指责；而你不做，还是会受到指责。”

\subsection*{原因之四：习惯}


我们的生活更多地反映了我们的习惯而不是我们所受到的教育。上学时在看过明星阿诺德·施瓦辛格主演的电影《野蛮人科南》之后，一个朋友说：“我想拥有像施瓦辛格那样的身材。”大部分男生都点头同意。

“但我听说他曾经很瘦弱。”另一个朋友补充说。

“是的，我也听说过，”另一位说道，“听说他几乎每天都泡在健身房里。”

“没错，我敢打赌他不得不这样。”

“不是的，”我们一伙人里总是好讽刺挖苦的那个人说，“我肯定他天生如此。算了吧，咱们别再谈论他了，去喝点啤酒吧。”

这是习惯控制行为的一个例子。我记得曾经问过富爸爸富人有什么习惯，他没有直接回答我，像往常一样他希望我从实例中学习。

“你爸爸什么时候支付账单？”富爸爸问道。

“月初。”我说。

“支付完账单之后他还有余钱吗？”他问。

“非常少。”我回答。

“这就是他苦苦挣扎的主要原因，”富爸爸说，“他有一些坏习惯。”

“你爸爸总是先支付给其他人，最后才支付给自己，而且这还得看他有无余钱。”

“他也并不是常没有余钱。”我说，“但他要按时支付账单，不是吗？难道你是说他不应该支付账单吗？”

“当然不是，”富爸爸说，“我坚信应该按时支付账单，只是我会先付给自己，再用剩余的钱给别人甚至政府。”

“但是如果你没有足够的钱，”我问，“又怎么办呢？”

“还是一样，”富爸爸说，“我仍然先支付给自己，即使我缺钱。因为对我来说，我的资产项比政府重要得多。”

“可是，”我说，“他们不会来找你要债吗？”

“会的，如果你不支付的话，”富爸爸说，“但是你看，我并没有说不支付。我只是说首先支付给自己，即便是我缺钱。”

“但是，”我又问，“你是怎么做的呢？”

“不是怎么做，而是为什么要这么做。”富爸爸说。

“那好，为什么？”

“这是动力问题，”富爸爸说，“如果我不支付给自己或是不支付给我的债主，你认为谁的抱怨声会更大些？”

“当然是你的债主会比你叫得响。”我说，这是显而易见的，“如果你不支付给自己的话，我想你也不会说什么。”

“所以你看，在我把仅有的钱先支付给自己后，要支付税款和其他债主的压力就会变得非常大，迫使我去寻求其他形式的收入，支付的压力就会成为我的动力。我会干额外的工作，开其他的公司，投资股票市场以及去做任何可以使那些人不再向我叫嚷的事。压力迫使我努力工作，迫使我去思考，最重要的是迫使我在钱的问题上更精明、更积极主动。如果我像你爸爸一样最后支付给自己，我虽然不会感到任何压力，但一定会因此破产。”

“你是说因为你欠了政府和其他人的债，所以受到了激励？”

“对，”富爸爸说，“你看，政府的征税部门和其他债主都有权有势，大部分人会屈服于这种威势，所以他们总是先支付这些账单而不去支付给自己。你听说过瘦弱的人被人把沙子踢到脸上的故事吧？”

我点了点头，“我在连环画里的举重和塑身广告里总是能看到这样的画面。”

“是的，大部分人让那些债主把沙子踢到脸上，而我决定利用对债主的恐惧使自己变得更加强壮。与此同时，其他人变得软弱可欺。我强迫自己考虑如何挣到额外的钱，就好比我去健身房做负重练习，我精神上的‘金钱肌肉’越发达，我就越强大。现在，我不再害怕那些人了。”

我喜欢富爸爸的这句话。“所以，如果我也学会先支付给自己，就会在财务上更‘强壮’，噢，应该是在精神上和财务上都更加‘强壮’。”

富爸爸点了点头。

“而如果我最后支付给自己，或根本就不支付，我就会变得越来越‘虚弱’，那么我一生都会围着老板、经理、税务官员、收账员和土地主们转，这仅仅是因为我没有良好的理财习惯。”

富爸爸点头称是，“就像故事里那个瘦弱的人一样。”

\subsection*{原因之五：傲慢}

傲慢是自大和无知的结合体。

“我知道的东西给我带来金钱，我不知道的东西使我失去金钱。因为每当我自高自大时，我就认为我不知道的东西并不重要。”富爸爸经常这样告诉我。

我发现许多人试图用傲慢来掩饰无知，当我同会计甚至是其他投资者讨论财务报告时，也经常会发生这样的事。

他们试图用自吹自擂来赢得争论，而我很清楚，这是因为他们根本不知道自己在说些什么。他们并没有撒谎，只是没有讲出真相。

在金钱、财务和投资领域，有许多人完全不知道自己在谈论什么。

财经领域的大部分人只不过像二手车推销员似的一心想兜售产品而已。

如果你知道自己在某一问题上有所欠缺，你就应该找一位本领域的专家或是一本相关的书，马上开始教育自己。

